% Manuale del Progettista - Progetto Ruggine
\documentclass[a4paper,12pt]{book}
\usepackage[utf8]{inputenc}
\usepackage[italian]{babel}
\usepackage{hyperref}
\usepackage{graphicx}

\title{Manuale del Progettista}
\author{Gruppo 50}
\date{\today}

\begin{document}
% Frontespizio
\begin{titlepage}
    \centering
    \includegraphics[width=0.3\textwidth]{img/ruggine_logo.png}\\
    {\Huge\bfseries Manuale del Progettista\par}
    \vspace{0.5cm}
    {\Large Progetto \textbf{Ruggine}\par}
    \vspace{0.5cm}
    {\large Corso di Programmazione di Sistema \\ 2025/2026\par}
    \vspace{1.5cm}
    {\large\textbf{Gruppo 50}}\\[0.5cm]
    Alessandro Annecchini \\ Riccardo Marconi \\ Francesco Palmisani \\ Giosuè Pinto
    \vfill
    {\large \today\par}
\end{titlepage}

\frontmatter
\tableofcontents

\mainmatter
\chapter{Introduzione}
Questo manuale descrive l'architettura, la logica e le modalità di sviluppo del progetto Ruggine.

\chapter{Panoramica dell'architettura}
\section{Struttura delle cartelle}
\section{Componenti principali (backend Rust, frontend React)}
\section{Flusso dati generale}

\chapter{Descrizione dei file principali}
\section{File Rust (main.rs, router.rs, ecc.)}
\section{Cartelle (controllers, services, repositories, dtos)}

\chapter{Descrizione delle funzioni}
Per ogni file, elenco e spiegazione delle funzioni principali e del loro ruolo.

\chapter{Flusso dei dati e logica applicativa}
\section{Autenticazione}
\section{Gestione messaggi e gruppi}
\section{Gestione inviti e notifiche}

\chapter{Estendibilità e personalizzazione}
\section{Aggiungere nuove funzionalità}
\section{Best practice per modifiche future}

\end{document}
